\documentclass[preprint,12pt]{elsarticle}

\usepackage[utf8]{inputenc}
\usepackage[T1]{fontenc}
\usepackage{amsmath,amssymb}
\usepackage{booktabs}
\usepackage{multirow}
\usepackage{caption}
\usepackage{float}
\usepackage{hyperref}
\usepackage{xcolor}
\usepackage{framed}

\definecolor{shadecolor}{RGB}{255,255,190}
\newenvironment{revision}{\begin{shaded}}{\end{shaded}}
\newcommand{\revphrase}[1]{\colorbox{yellow!30}{\strut #1}}
\captionsetup{justification=centering}

\makeatletter
\long\def\pprintMaketitle{\clearpage
  \resetTitleCounters
  \def\baselinestretch{1}%
  \printFirstPageNotes
  \begin{\elsarticletitlealign}%
    \thispagestyle{pprintTitle}%
    \def\baselinestretch{1}%
    \Large\@title\par\vskip6pt
  \end{\elsarticletitlealign}%
  \gdef\thefootnote{\arabic{footnote}}%
}
\makeatother

\renewcommand{\thetable}{S\arabic{table}}
\renewcommand{\thesection}{S\arabic{section}}

\journal{Computer Methods and Programs in Biomedicine}

\begin{document}

\begin{frontmatter}
\title{Supplementary Material\\[15pt]
\large Windowed Variational Mode Decomposition with Gaussian-Kernel Mode Selection for Sex-Stratified Speech-Based Depression Screening}
\vspace{-1cm}

\end{frontmatter}

\section{Window-Duration Comparison (0.5\,s, 1\,s, 2\,s)}
\label{sec:supp-window}

This supplement compares the three window durations evaluated in the window-duration analysis
(Section 3.6 of the main text): 0.5\,s, 1\,s, and 2\,s, for both sexes,
under the same discontinuity-aware pipeline reported throughout the main text. Within each model
category, architecture selection was performed from out-of-fold TRAIN-CV AUC-ROC without using
TEST (checkpoint 6 of Section 2.12). Table~\ref{tab:supp-window-best} reports the resulting
category representative used for each sex--window configuration after Youden threshold
calibration, together with its fixed-split TEST evaluation (Accuracy, F1) and its Experiment 3a
mean AUC-ROC (fixed split, varied model seeds, $N_{\text{reps}}=10$). These TEST quantities are
reported descriptively and should not be interpreted as an independent estimate after all design
choices.

\begin{table}[H]
\centering
\captionsetup{justification=raggedright, singlelinecheck=false}
\caption{Reported architecture by window duration and sex. Within each category,
the representative was selected from TRAIN-CV (checkpoint 6, Section 2.12); among the three
category representatives, the row reports the one with the highest calibrated F1 on the original
fixed TEST split. AUC-ROC is the Experiment 3a mean ($N_{\text{reps}}=10$); F1/Accuracy are
descriptive values from the calibrated TEST evaluation.}
\label{tab:supp-window-best}
\resizebox{\linewidth}{!}{%
\begin{tabular}{llp{5.4cm}ccc}
\toprule
Window & Sex & Reported architecture & AUC-ROC (3a) & F1 & Accuracy \\
\midrule
\multirow{2}{*}{0.5\,s}
 & FEMALE & Individual: DT & 0.529 & 0.533 & 0.632 \\[8pt]
 & MALE   & Individual: LR & \normalsize\textbf{0.594} & \normalsize\textbf{0.421} & \normalsize\textbf{0.500} \\
\midrule
\multirow{2}{*}{1\,s}
 & FEMALE & Stacking-2: \newline GB + XGBoost $\rightarrow$ RF &  \normalsize\textbf{0.617} & \normalsize\textbf{0.526} & \normalsize\textbf{0.526} \\\addlinespace[8pt]
 & MALE   & Stacking-2: \newline DT + RF $\rightarrow$ XGBoost & 0.459 & 0.526 & 0.591 \\
\midrule
\multirow{2}{*}{2\,s}
 & FEMALE & Stacking-2: \newline GB + XGBoost $\rightarrow$ DT & 0.612 & 0.429 & 0.579 \\\addlinespace[8pt]
 & MALE   & Stacking-2: \newline DT + RF $\rightarrow$ XGBoost & 0.556 & 0.571 & 0.727 \\
\bottomrule
\end{tabular}
}
\par\vspace{0.5em}
{\small\raggedright\emph{Note:} DT $=$ Decision Tree; RF $=$ Random Forest; GB $=$ Gradient Boosting.
Within each architecture category, the representative was selected using TRAIN-CV AUC-ROC rather
than TEST. The single architecture reported for each sex$\times$window configuration is the
category representative with the highest calibrated F1 on the original fixed TEST split. This
cross-category TEST-F1 comparison is descriptive, not an independent TRAIN-only selection step.
Bold: each sex's headline window (FEMALE 1\,s, MALE 0.5\,s).
Table~\ref{tab:supp-window-wc} reports the corresponding repeated-split sensitivity analysis
(Experiment 3b).\par}
\end{table}

Table~\ref{tab:supp-window-wc} reports Experiment 3b, a repeated-split sensitivity analysis with
randomized TRAIN/TEST partitions ($N_{\text{reps}}=10$; Section 2.11 of the main text). Because
the repeated partitions are drawn from the same finite subject pool and are not independent, the
reported mean $\pm 1.96\,\mathrm{SE}$ intervals characterize variation across the evaluated
splits and are not used here as formal significance tests against AUC-ROC $=0.5$. Among the six
sex-by-window combinations, FEMALE at 1\,s and MALE at 0.5\,s (bold below) are the two for which
the repeated-split mean remained above 0.5 and the reported resampling interval did not span 0.5.
These are treated as the headline configurations in the main text, while all six combinations are
reported here for completeness. This choice remains exploratory because window duration itself was
not selected inside a fully nested resampling procedure.

\begin{table}[H]
\centering
\begin{minipage}{0.75\linewidth}
\captionsetup{justification=raggedright, singlelinecheck=false}
\caption{Repeated-split sensitivity analysis by window duration and sex: AUC-ROC on randomized TRAIN/TEST splits (Experiment 3b, $N_{\text{reps}}=10$), using the corresponding TRAIN-selected within-category architecture reported in Table~\ref{tab:supp-window-best}.}
\label{tab:supp-window-wc}
\resizebox{\linewidth}{!}{%
\begin{tabular}{llccc}
\toprule
Window & Sex & Mean AUC & Mean $\pm1.96\,\mathrm{SE}$ interval & Interval spans 0.5 \\
\midrule
\multirow{2}{*}{0.5\,s}
 & FEMALE & 0.471 & [0.420-0.523] & Yes \\\addlinespace[4pt]
 & MALE   & \fontsize{11}{13}\selectfont\textbf{0.604} & \fontsize{10}{12}\selectfont\textbf{[0.537-0.672]} & \fontsize{10}{12}\selectfont\textbf{No} \\
\midrule
\multirow{2}{*}{1\,s}
 & FEMALE & \fontsize{11}{13}\selectfont\textbf{0.611} & \fontsize{10}{12}\selectfont\textbf{[0.549-0.674]} & \fontsize{10}{12}\selectfont\textbf{No} \\\addlinespace[4pt]
 & MALE   & 0.511 & [0.370-0.653] & Yes \\
\midrule
\multirow{2}{*}{2\,s}
 & FEMALE & 0.488 & [0.395-0.581] & Yes \\\addlinespace[4pt]
 & MALE   & 0.571 & [0.488-0.654] & Yes \\
\bottomrule
\end{tabular}
}
\par\vspace{0.5em}
{\small\raggedright\emph{Note:} Bold $=$ each sex's headline window (FEMALE 1\,s, MALE 0.5\,s). The intervals are mean $\pm1.96\,\mathrm{SE}$ across 10 repeated splits of the same finite subject pool; they are descriptive resampling intervals, not formal independent-sample 95\% confidence intervals.\par}
\end{minipage}
\end{table}

The fixed-split results in Table~\ref{tab:supp-window-best} illustrate why conclusions based on a
single partition would be unstable. For example, FEMALE at 2\,s has a fixed-split AUC-ROC of
0.612, whereas its Experiment 3b repeated-split mean is 0.488 with a mean $\pm1.96\,\mathrm{SE}$
interval of [0.395--0.581]. The other non-headline combinations likewise show substantial
partition sensitivity. FEMALE at 1\,s and MALE at 0.5\,s are the two configurations for which
the repeated-split means are above 0.5 and the corresponding descriptive intervals do not span
0.5 (0.611 [0.549--0.674] and 0.604 [0.537--0.672], respectively). This pattern motivated their
use as the headline configurations for descriptive reporting in Section 3.6 of the main text
rather than adopting a single sex-agnostic window length.
However, because window duration was itself compared empirically and was not selected within a
fully nested procedure, this should be interpreted as an exploratory robustness result rather
than confirmatory evidence that discrimination exists only at those two window lengths.

\section{Individual Classifier Performance at Non-Headline Windows}
\label{sec:supp-individual}

For completeness, Tables~\ref{tab:supp-individual-female-500ms}--\ref{tab:supp-individual-male-2s}
report the fixed-split TEST performance (default probability threshold, 0.5) of the eight base
classifiers described in Section 2.8 of the main text, at all three window durations for both
sexes, including each sex's own headline window (FEMALE at 1\,s, MALE at 0.5\,s). Main-text
Tables 5 and 6 report the corresponding headline-window individual-model results using
out-of-fold TRAIN-CV metrics for within-category ranking, whereas
Tables~\ref{tab:supp-individual-female-1s} and~\ref{tab:supp-individual-male-500ms} report
fixed-split TEST performance for the same eight classifiers, with bootstrap uncertainty
intervals. The two sets of tables therefore serve different purposes and are not expected to have
the same ordering or numerical values. Values in brackets are the reported 95\%
percentile-bootstrap intervals over TEST subjects. Given the small TEST samples in each sex,
these intervals are expected to be wide and are presented as uncertainty summaries rather than
evidence for fine-grained differences between classifiers. Rows are sorted by TEST AUC-ROC for
descriptive completeness only. This ordering is not the architecture-selection criterion used
within the individual-model category, which instead relies on out-of-fold TRAIN-CV AUC-ROC
without consulting TEST for the within-category ranking.

\begin{table}[H]
\centering
\captionsetup{justification=raggedright, singlelinecheck=false}
\caption{Performance of the eight individual classifiers on TEST, 0.5\,s window, FEMALE (default
threshold = 0.5). Value [95\% \emph{bootstrap} CI].}
\label{tab:supp-individual-female-500ms}
\resizebox{\linewidth}{!}{%
\renewcommand{\arraystretch}{1.6}
\begin{tabular}{lcccccc}
\toprule
{\large Model} & {\large AUC-ROC} & {\large Accuracy} & {\large Sensitivity} & {\large Precision} & {\large Specificity} & {\large F1} \\
\midrule
SVM                  & 0.810 [0.550--1.000] & 0.263 [0.053--0.474] & 0.286 [0.000--0.667] & 0.182 [0.000--0.444] & 0.250 [0.000--0.500] & 0.222 [0.000--0.476] \\
Random Forest        & 0.702 [0.389--0.957] & 0.632 [0.421--0.842] & 0.143 [0.000--0.500] & 0.500 [0.000--1.000] & 0.917 [0.727--1.000] & 0.222 [0.000--0.571] \\
XGBoost              & 0.690 [0.405--0.933] & 0.632 [0.421--0.842] & 0.286 [0.000--0.667] & 0.500 [0.000--1.000] & 0.833 [0.600--1.000] & 0.364 [0.000--0.667] \\
Decision Tree        & 0.619 [0.369--0.847] & 0.632 [0.421--0.842] & 0.571 [0.167--1.000] & 0.500 [0.143--0.857] & 0.667 [0.375--0.917] & 0.533 [0.167--0.800] \\
Gradient Boosting    & 0.500 [0.148--0.829] & 0.737 [0.526--0.895] & 0.429 [0.000--0.833] & 0.750 [0.000--1.000] & 0.917 [0.727--1.000] & 0.545 [0.000--0.857] \\
Logistic Regression  & 0.310 [0.023--0.643] & 0.421 [0.211--0.632] & 0.286 [0.000--0.667] & 0.250 [0.000--0.600] & 0.500 [0.214--0.778] & 0.267 [0.000--0.556] \\
Naive Bayes          & 0.262 [0.000--0.590] & 0.316 [0.105--0.526] & 0.286 [0.000--0.667] & 0.200 [0.000--0.500] & 0.333 [0.083--0.615] & 0.235 [0.000--0.500] \\
KNN                  & 0.196 [0.038--0.405] & 0.368 [0.158--0.579] & 0.000 [0.000--0.000] & 0.000 [0.000--0.000] & 0.583 [0.300--0.857] & 0.000 [0.000--0.000] \\
\bottomrule
\end{tabular}%
}
\end{table}

\begin{table}[H]
\centering
\captionsetup{justification=raggedright, singlelinecheck=false}
\caption{Performance of the eight individual classifiers on TEST, 1\,s window (FEMALE headline
window), FEMALE (default threshold = 0.5). Value [95\% \emph{bootstrap} CI]. TEST-set companion
to main-text Table~5, which reports out-of-fold TRAIN-CV metrics for within-category ranking.}
\label{tab:supp-individual-female-1s}
\resizebox{\linewidth}{!}{%
\renewcommand{\arraystretch}{1.6}
\begin{tabular}{lcccccc}
\toprule
{\large Model} & {\large AUC-ROC} & {\large Accuracy} & {\large Sensitivity} & {\large Precision} & {\large Specificity} & {\large F1} \\
\midrule
Logistic Regression  & 0.738 [0.461--0.952] & 0.684 [0.474--0.895] & 0.714 [0.333--1.000] & 0.556 [0.200--0.889] & 0.667 [0.385--0.917] & 0.625 [0.286--0.857] \\
SVM                  & 0.643 [0.318--0.920] & 0.632 [0.421--0.842] & 0.000 [0.000--0.000] & 0.000 [0.000--0.000] & 1.000 [1.000--1.000] & 0.000 [0.000--0.000] \\
Naive Bayes          & 0.548 [0.318--0.750] & 0.526 [0.316--0.737] & 1.000 [1.000--1.000] & 0.438 [0.200--0.688] & 0.250 [0.000--0.500] & 0.609 [0.333--0.815] \\
Random Forest        & 0.512 [0.189--0.821] & 0.579 [0.368--0.789] & 0.286 [0.000--0.667] & 0.400 [0.000--1.000] & 0.750 [0.500--1.000] & 0.333 [0.000--0.667] \\
KNN                  & 0.500 [0.196--0.800] & 0.632 [0.421--0.842] & 0.429 [0.000--0.833] & 0.500 [0.000--1.000] & 0.750 [0.500--1.000] & 0.462 [0.000--0.750] \\
Gradient Boosting    & 0.464 [0.171--0.779] & 0.368 [0.158--0.579] & 0.714 [0.333--1.000] & 0.333 [0.111--0.583] & 0.167 [0.000--0.400] & 0.455 [0.190--0.692] \\
Decision Tree        & 0.452 [0.211--0.692] & 0.421 [0.211--0.632] & 0.571 [0.167--1.000] & 0.333 [0.083--0.615] & 0.333 [0.083--0.615] & 0.421 [0.125--0.667] \\
XGBoost              & 0.452 [0.170--0.750] & 0.526 [0.316--0.737] & 0.571 [0.167--1.000] & 0.400 [0.100--0.714] & 0.500 [0.214--0.786] & 0.471 [0.143--0.737] \\
\bottomrule
\end{tabular}%
}
\end{table}

\begin{table}[H]
\centering
\captionsetup{justification=raggedright, singlelinecheck=false}
\caption{Performance of the eight individual classifiers on TEST, 2\,s window, FEMALE (default
threshold = 0.5). Value [95\% \emph{bootstrap} CI].}
\label{tab:supp-individual-female-2s}
\resizebox{\linewidth}{!}{%
\renewcommand{\arraystretch}{1.6}
\begin{tabular}{lcccccc}
\toprule
{\large Model} & {\large AUC-ROC} & {\large Accuracy} & {\large Sensitivity} & {\large Precision} & {\large Specificity} & {\large F1} \\
\midrule
Gradient Boosting    & 0.690 [0.372--0.966] & 0.632 [0.421--0.842] & 0.714 [0.333--1.000] & 0.500 [0.182--0.818] & 0.583 [0.300--0.857] & 0.588 [0.250--0.833] \\
Random Forest        & 0.595 [0.311--0.864] & 0.526 [0.316--0.737] & 0.429 [0.000--0.833] & 0.375 [0.000--0.750] & 0.583 [0.300--0.857] & 0.400 [0.000--0.667] \\
XGBoost              & 0.560 [0.250--0.833] & 0.421 [0.211--0.632] & 0.429 [0.000--0.833] & 0.300 [0.000--0.600] & 0.417 [0.143--0.714] & 0.353 [0.000--0.625] \\
KNN                  & 0.500 [0.170--0.811] & 0.684 [0.474--0.895] & 0.143 [0.000--0.500] & 1.000 [0.000--1.000] & 1.000 [1.000--1.000] & 0.250 [0.000--0.667] \\
Decision Tree        & 0.470 [0.227--0.733] & 0.632 [0.421--0.842] & 0.286 [0.000--0.667] & 0.500 [0.000--1.000] & 0.833 [0.600--1.000] & 0.364 [0.000--0.667] \\
SVM                  & 0.440 [0.136--0.756] & 0.474 [0.263--0.684] & 0.429 [0.000--0.833] & 0.333 [0.000--0.667] & 0.500 [0.222--0.786] & 0.375 [0.000--0.667] \\
Logistic Regression  & 0.440 [0.143--0.744] & 0.474 [0.263--0.684] & 0.143 [0.000--0.500] & 0.200 [0.000--0.667] & 0.667 [0.375--0.917] & 0.167 [0.000--0.462] \\
Naive Bayes          & 0.393 [0.125--0.689] & 0.526 [0.316--0.737] & 0.429 [0.000--0.833] & 0.375 [0.000--0.750] & 0.583 [0.300--0.857] & 0.400 [0.000--0.667] \\
\bottomrule
\end{tabular}%
}
\end{table}

\begin{table}[H]
\centering
\captionsetup{justification=raggedright, singlelinecheck=false}
\caption{Performance of the eight individual classifiers on TEST, 0.5\,s window (MALE headline
window), MALE (default threshold = 0.5). Value [95\% \emph{bootstrap} CI]. TEST-set companion to
main-text Table~6, which reports out-of-fold TRAIN-CV metrics for within-category ranking.}
\label{tab:supp-individual-male-500ms}
\resizebox{\linewidth}{!}{%
\renewcommand{\arraystretch}{1.6}
\begin{tabular}{lcccccc}
\toprule
{\large Model} & {\large AUC-ROC} & {\large Accuracy} & {\large Sensitivity} & {\large Precision} & {\large Specificity} & {\large F1} \\
\midrule
XGBoost              & 0.677 [0.354--0.944] & 0.727 [0.545--0.909] & 0.333 [0.000--0.750] & 0.500 [0.000--1.000] & 0.875 [0.692--1.000] & 0.400 [0.000--0.750] \\
KNN                  & 0.651 [0.430--0.859] & 0.545 [0.318--0.739] & 0.000 [0.000--0.000] & 0.000 [0.000--0.000] & 0.750 [0.529--0.941] & 0.000 [0.000--0.000] \\
Random Forest        & 0.635 [0.362--0.882] & 0.773 [0.591--0.909] & 0.167 [0.000--0.500] & 1.000 [0.000--1.000] & 1.000 [1.000--1.000] & 0.286 [0.000--0.667] \\
Logistic Regression  & 0.594 [0.286--0.875] & 0.545 [0.318--0.727] & 0.667 [0.250--1.000] & 0.333 [0.083--0.615] & 0.500 [0.250--0.750] & 0.444 [0.133--0.700] \\
Naive Bayes          & 0.583 [0.233--0.906] & 0.273 [0.091--0.455] & 0.667 [0.250--1.000] & 0.222 [0.053--0.429] & 0.125 [0.000--0.312] & 0.333 [0.091--0.560] \\
SVM                  & 0.542 [0.219--0.859] & 0.500 [0.273--0.727] & 0.500 [0.000--1.000] & 0.273 [0.000--0.556] & 0.500 [0.250--0.750] & 0.353 [0.000--0.625] \\
Decision Tree        & 0.479 [0.263--0.728] & 0.545 [0.318--0.727] & 0.333 [0.000--0.750] & 0.250 [0.000--0.600] & 0.625 [0.375--0.857] & 0.286 [0.000--0.571] \\
Gradient Boosting    & 0.469 [0.135--0.807] & 0.591 [0.364--0.773] & 0.500 [0.000--1.000] & 0.333 [0.000--0.667] & 0.625 [0.375--0.857] & 0.400 [0.000--0.667] \\
\bottomrule
\end{tabular}%
}
\end{table}

\begin{table}[H]
\centering
\captionsetup{justification=raggedright, singlelinecheck=false}
\caption{Performance of the eight individual classifiers on TEST, 1\,s window, MALE (default
threshold = 0.5). Value [95\% \emph{bootstrap} CI].}
\label{tab:supp-individual-male-1s}
\resizebox{\linewidth}{!}{%
\renewcommand{\arraystretch}{1.6}
\begin{tabular}{lcccccc}
\toprule
{\large Model} & {\large AUC-ROC} & {\large Accuracy} & {\large Sensitivity} & {\large Precision} & {\large Specificity} & {\large F1} \\
\midrule
XGBoost              & 0.688 [0.376--0.950] & 0.727 [0.545--0.909] & 0.167 [0.000--0.500] & 0.500 [0.000--1.000] & 0.938 [0.789--1.000] & 0.250 [0.000--0.667] \\
Gradient Boosting    & 0.677 [0.350--0.958] & 0.727 [0.545--0.909] & 0.667 [0.217--1.000] & 0.500 [0.143--0.857] & 0.750 [0.529--0.941] & 0.571 [0.182--0.842] \\
Decision Tree        & 0.656 [0.395--0.895] & 0.682 [0.500--0.864] & 0.500 [0.000--1.000] & 0.429 [0.000--0.833] & 0.750 [0.529--0.941] & 0.462 [0.000--0.750] \\
SVM                  & 0.573 [0.188--0.930] & 0.727 [0.545--0.909] & 0.000 [0.000--0.000] & 0.000 [0.000--0.000] & 1.000 [1.000--1.000] & 0.000 [0.000--0.000] \\
Random Forest        & 0.526 [0.224--0.819] & 0.727 [0.545--0.909] & 0.000 [0.000--0.000] & 0.000 [0.000--0.000] & 1.000 [1.000--1.000] & 0.000 [0.000--0.000] \\
Naive Bayes          & 0.500 [0.209--0.789] & 0.727 [0.545--0.909] & 0.167 [0.000--0.500] & 0.500 [0.000--1.000] & 0.938 [0.800--1.000] & 0.250 [0.000--0.667] \\
Logistic Regression  & 0.448 [0.105--0.810] & 0.500 [0.318--0.682] & 0.333 [0.000--0.778] & 0.222 [0.000--0.545] & 0.562 [0.316--0.800] & 0.267 [0.000--0.545] \\
KNN                  & 0.448 [0.194--0.714] & 0.727 [0.545--0.909] & 0.000 [0.000--0.000] & 0.000 [0.000--0.000] & 1.000 [1.000--1.000] & 0.000 [0.000--0.000] \\
\bottomrule
\end{tabular}%
}
\end{table}

\begin{table}[H]
\centering
\captionsetup{justification=raggedright, singlelinecheck=false}
\caption{Performance of the eight individual classifiers on TEST, 2\,s window, MALE (default
threshold = 0.5). Value [95\% \emph{bootstrap} CI].}
\label{tab:supp-individual-male-2s}
\resizebox{\linewidth}{!}{%
\renewcommand{\arraystretch}{1.6}
\begin{tabular}{lcccccc}
\toprule
{\large Model} & {\large AUC-ROC} & {\large Accuracy} & {\large Sensitivity} & {\large Precision} & {\large Specificity} & {\large F1} \\
\midrule
SVM                  & 0.625 [0.316--0.912] & 0.727 [0.545--0.909] & 0.000 [0.000--0.000] & 0.000 [0.000--0.000] & 1.000 [1.000--1.000] & 0.000 [0.000--0.000] \\
Naive Bayes          & 0.562 [0.292--0.821] & 0.591 [0.364--0.773] & 0.667 [0.222--1.000] & 0.364 [0.091--0.667] & 0.562 [0.316--0.800] & 0.471 [0.143--0.727] \\
Logistic Regression  & 0.500 [0.500--0.500] & 0.727 [0.545--0.909] & 0.000 [0.000--0.000] & 0.000 [0.000--0.000] & 1.000 [1.000--1.000] & 0.000 [0.000--0.000] \\
Decision Tree        & 0.365 [0.167--0.639] & 0.409 [0.227--0.636] & 0.167 [0.000--0.500] & 0.111 [0.000--0.364] & 0.500 [0.250--0.750] & 0.133 [0.000--0.375] \\
XGBoost              & 0.333 [0.106--0.597] & 0.500 [0.273--0.727] & 0.000 [0.000--0.000] & 0.000 [0.000--0.000] & 0.688 [0.450--0.895] & 0.000 [0.000--0.000] \\
KNN                  & 0.266 [0.028--0.580] & 0.682 [0.500--0.864] & 0.000 [0.000--0.000] & 0.000 [0.000--0.000] & 0.938 [0.789--1.000] & 0.000 [0.000--0.000] \\
Random Forest        & 0.229 [0.053--0.444] & 0.682 [0.500--0.864] & 0.000 [0.000--0.000] & 0.000 [0.000--0.000] & 0.938 [0.789--1.000] & 0.000 [0.000--0.000] \\
Gradient Boosting    & 0.219 [0.048--0.447] & 0.455 [0.273--0.636] & 0.000 [0.000--0.000] & 0.000 [0.000--0.000] & 0.625 [0.385--0.850] & 0.000 [0.000--0.000] \\
\bottomrule
\end{tabular}%
}
\end{table}

Consistent with the pattern already reported for each sex's headline window (Section 3.2 of the
main text), no single classifier reaches strong, consistent performance across both sexes at any
of the non-headline window durations: several classifiers collapse to predicting only the
majority class (HEALTH) at the default threshold regardless of window length, most visibly for
MALE at 2\,s, where five of the eight models never predict a single DEP case. These per-model,
default-threshold TEST rankings are, in every window and sex shown here, distinct from the
architecture selected within category by out-of-fold TRAIN-CV AUC-ROC and reported in
Table~\ref{tab:supp-window-best}. The tables are therefore descriptive TEST-set summaries and
should not be read as an alternative TEST-based model-selection procedure.

\section{Partially Nested Cross-Validation: Mode-Ranking Stability and Preprocessing-Bias Diagnostic}
\label{sec:supp-nested}

As noted in Section 2.12 of the main text, the Gaussian-kernel mode ranking and the global scaler
are each fit once on all TRAIN subjects before the 5-fold cross-validation used to tune
classifier hyperparameters, so a held-out subject's own label and window distribution have
already contributed to the mode ranking and normalization applied to their own fold. This does
not introduce information from TEST into TRAIN, but it can make the internal TRAIN-CV scores
optimistic. To examine the magnitude of this specific preprocessing-related source of bias, a
partially nested cross-validation was run for all six sex-by-window combinations: five stratified
outer folds were built over the TRAIN subjects, and within each fold the mode ranking, scaler,
and each classifier's hyperparameters were refit using only that fold's own training subjects.
Window length, $K$, and $N_{\text{select}}$ were held fixed at their reported values. Therefore,
this diagnostic addresses refitting of mode ranking, scaling and classifier tuning conditional on
those upstream design choices; it does not quantify the additional selection bias associated with
choosing window length, $K$, or $N_{\text{select}}$.

Table~\ref{tab:supp-nested-modes} reports the Gaussian-kernel mode set obtained independently in
each of the five outer folds. Several configurations are internally stable across folds, but
fold-to-fold stability should not be interpreted as validation of the mode set used in the
reported pipeline. In particular, at MALE 0.5\,s all five outer folds select
$\{0,1,2,6,7\}$, whereas the reported MALE configuration uses
$\{0,1,2,3,7\}$. Likewise, at FEMALE 0.5\,s all five folds select
$\{0,1,2,7\}$ rather than the reported FEMALE set $\{0,1,2,3\}$.
At FEMALE 1\,s, four of the five folds agree with the reported set
$\{0,1,2,3\}$, with one fold selecting $\{0,2,3,7\}$. At MALE 1\,s,
all five folds select the reported set $\{0,1,2,3,7\}$. At 2\,s,
FEMALE shows three distinct fold-specific sets, whereas MALE is internally stable at
$\{0,1,2,3,4\}$. Thus, Table~\ref{tab:supp-nested-modes} informs both fold-to-fold
reproducibility and agreement with the reported mode choice, which are separate properties.

\begin{table}[H]
\centering
\begin{minipage}{0.8\linewidth}
\captionsetup{justification=raggedright, singlelinecheck=false}
\caption{Gaussian-kernel mode set obtained independently in each of five stratified outer folds of
the partially nested cross-validation, by window duration and sex. The ``Stable'' column refers
only to agreement across the five outer folds, not to agreement with the mode set used in the
reported pipeline.}
\label{tab:supp-nested-modes}
\resizebox{\linewidth}{!}{%
\begin{tabular}{llp{7.5cm}c}
\toprule
Window & Sex & Modes selected per outer fold (fold 1--5) & Stable \\
\midrule
\multirow{2}{*}{0.5\,s}
 & FEMALE & \{0,1,2,7\} in every fold & Yes \\\addlinespace[8pt]
 & MALE   & \{0,1,2,6,7\} in every fold & Yes \\
\midrule
\multirow{2}{*}{1\,s}
 & FEMALE & \{0,1,2,3\}; \{0,1,2,3\}; \{0,1,2,3\}; \{0,2,3,7\}; \{0,1,2,3\} & 4/5 \\\addlinespace[8pt]
 & MALE   & \{0,1,2,3,7\} in every fold & Yes \\
\midrule
\multirow{2}{*}{2\,s}
 & FEMALE & \{0,2,3,4\}; \{0,1,3,4\}; \{0,2,3,4\}; \{0,1,3,4\}; \{1,2,3,4\} & No \\\addlinespace[8pt]
 & MALE   & \{0,1,2,3,4\} in every fold & Yes \\
\bottomrule
\end{tabular}%
}
\end{minipage}
\end{table}

Table~\ref{tab:supp-nested-auc} reports the resulting partially nested AUC-ROC (selected model per
outer fold, mean $\pm$ SD across folds), alongside Experiment 3b, the repeated-split sensitivity
analysis in Table~\ref{tab:supp-window-wc}, for the same window and sex. The two sets of values
are not directly comparable because they use different resampling schemes and answer different
questions. The five-fold partially nested analysis is included only as a diagnostic of whether
refitting mode ranking, scaling and classifier tuning within the outer folds materially changes
the observed performance pattern, conditional on fixed window length, $K$, and
$N_{\text{select}}$; it is not an additional headline estimate or a correction for all sources
of model-selection bias.

\begin{table}[H]
\centering
\captionsetup{justification=raggedright, singlelinecheck=false}
\caption{Partially nested AUC-ROC (selected model per outer fold, mean $\pm$ SD across the
five folds) versus the Experiment 3b repeated-split sensitivity analysis
(Table~\ref{tab:supp-window-wc}), by window duration and sex.}
\label{tab:supp-nested-auc}
\resizebox{\linewidth}{!}{%
\begin{tabular}{llccc}
\toprule
Window & Sex & \begin{tabular}[t]{@{}c@{}}Nested-CV AUC-ROC\\(mean $\pm$ SD)\end{tabular} & Nested-CV range & \begin{tabular}[t]{@{}c@{}}Exp.\ 3b AUC-ROC\\{[}mean $\pm1.96\,\mathrm{SE}$ interval{]}\end{tabular} \\
\midrule
\multirow{2}{*}{0.5\,s}
 & FEMALE & 0.677 $\pm$ 0.177 & [0.472--0.900] & 0.471 [0.420--0.523] \\\addlinespace[10pt]
 & MALE   & 0.716 $\pm$ 0.128 & [0.583--0.896] & 0.604 [0.537--0.672] \\
\midrule
\multirow{2}{*}{1\,s}
 & FEMALE & 0.612 $\pm$ 0.136 & [0.472--0.800] & 0.611 [0.549--0.674] \\\addlinespace[10pt]
 & MALE   & 0.626 $\pm$ 0.110 & [0.542--0.818] & 0.511 [0.370--0.653] \\
\midrule
\multirow{2}{*}{2\,s}
 & FEMALE & 0.670 $\pm$ 0.091 & [0.528--0.775] & 0.488 [0.395--0.581] \\\addlinespace[10pt]
 & MALE   & 0.674 $\pm$ 0.134 & [0.455--0.792] & 0.571 [0.488--0.654] \\
\bottomrule
\end{tabular}%
}
\par\vspace{0.5em}
{\small\raggedright\emph{Note:} Nested-CV AUC-ROC uses five outer folds, whereas Experiment 3b uses 10 repeated TRAIN/TEST splits of the same finite subject pool. The two columns summarize different resampling procedures and are shown only as a sensitivity comparison; numerical agreement is not expected.\par}
\end{table}

Across all six sex-by-window combinations, the partially nested AUC-ROC values are of broadly
similar magnitude to the corresponding Experiment 3b means, although the two procedures are not
directly comparable. FEMALE at 1\,s shows the closest numerical agreement (0.612 in the partially
nested analysis versus 0.611 in Experiment 3b). However, the mode-set results in
Table~\ref{tab:supp-nested-modes} show that internal fold stability does not necessarily imply
agreement with the mode set used in the reported pipeline. This is particularly clear at
MALE 0.5\,s and FEMALE 0.5\,s, where all five outer folds agree with each other but select a
different set from the one reported. The partially nested analysis should therefore be interpreted
as a diagnostic of sensitivity to refitting within folds, not as validation of the original mode
subset. It also does not rule out optimism arising from upstream choices held fixed in this
diagnostic, including window duration, $K$, $N_{\text{select}}$, or other researcher degrees of
freedom.

\end{document}
